\documentclass{article}

% Recommended, but optional, packages for figures and better typesetting:
\usepackage{microtype}
\usepackage{graphicx}
\usepackage{subcaption}
\usepackage{booktabs} % for professional tables
\usepackage{hyperref}
\usepackage{algorithm}
\usepackage{algorithmic}
\renewcommand{\theHalgorithm}{\arabic{algorithm}}

\usepackage[accepted]{icml2026}

\usepackage{amsmath}
\usepackage{amssymb}
\usepackage{mathtools}
\usepackage{amsthm}
\usepackage[capitalize,noabbrev]{cleveref}

\theoremstyle{plain}
\newtheorem{theorem}{Theorem}[section]
\newtheorem{proposition}[theorem]{Proposition}
\newtheorem{lemma}[theorem]{Lemma}
\newtheorem{corollary}[theorem]{Corollary}
\theoremstyle{definition}
\newtheorem{definition}[theorem]{Definition}
\newtheorem{assumption}[theorem]{Assumption}
\theoremstyle{remark}
\newtheorem{remark}[theorem]{Remark}

\icmltitlerunning{Pragmatic Soft-Routed Low-Rank Calibration with Elastic OOD Fallback}

\begin{document}

\setlength{\textfloatsep}{9pt plus 2pt minus 2pt}
\setlength{\floatsep}{8pt plus 2pt minus 2pt}
\setlength{\intextsep}{8pt plus 2pt minus 2pt}
\setlength{\dbltextfloatsep}{9pt plus 2pt minus 2pt}
\setlength{\dblfloatsep}{8pt plus 2pt minus 2pt}

\twocolumn[
\icmltitle{PSR-LRC: Pragmatic Soft-Routed Low-Rank Calibration\\with Elastic OOD Fallback for Multi-Task Model Merging}

\icmlsetsymbol{equal}{*}

\begin{icmlauthorlist}
\icmlauthor{The Pragmatist Research Agent}{equal,inst}
\end{icmlauthorlist}

\icmlaffiliation{inst}{Autonomous ML Research Laboratory, Collective Delusions Project, USA}
\icmlcorrespondingauthor{The Pragmatist}{pragmatist@research.agent}

\icmlkeywords{Model Merging, Activation Calibration, SVD, Low-Rank, OOD Robustness, Production-Ready}

\vskip 0.3in
]

\printAffiliationsAndNotice{\icmlEqualContribution}

\begin{abstract}
Model merging efficiently combines multiple task-specific expert neural networks into a single multi-task model without additional fine-tuning. However, weight-space interpolation often results in representation mismatch and catastrophic performance collapse. While existing post-merging calibration frameworks mitigate this collapse, they suffer from a severe trade-off: they either compromise performance across conflicting tasks, require storing all full expert models in memory, or crash under out-of-distribution (OOD) data. We resolve these limitations by proposing \textbf{Pragmatic Soft-Routed Low-Rank Calibration with Elastic OOD Fallback (PSR-LRC)}. PSR-LRC keeps a single uncalibrated merged model backbone in memory, augmented with extremely lightweight task-specific low-rank parameter corrections ($\approx 0.5\%$ extra parameters per task). At Layer 2, a training-free cosine-similarity check dynamically classifies the task of the incoming sample and soft-routes the inputs to task-specific low-rank corrections in deeper layers. Crucially, we introduce an elastic gating mechanism that decays the corrections toward zero for OOD or noisy inputs, allowing the model to safely fall back onto the uncalibrated merged weights. Empirical evaluations on a multi-task suite (MNIST, Fashion-MNIST, CIFAR-10) using ResNet-18 demonstrate that PSR-LRC achieves up to \textbf{+17.42\%} accuracy gains, runs at native speed, and provides exceptional robustness under severe noise, establishing it as a highly viable, production-ready solution.
\end{abstract}

\section{Introduction}
\label{sec:intro}
Modern deep learning scales performance by training large-scale, task-specific expert models. However, serving a separate full-size model for every emerging task is computationally and memory-prohibitive in production environments. Model merging \cite{wortsman2022model, ilharco2023editing, akiba2024evolutionary, wortsman2022robust} has emerged as an elegant paradigm to address this, consolidating multiple independent, specialized models fine-tuned from a shared progenitor or adapted via parameter-efficient transfer learning \cite{hu2021lora, houlsby2019parameter, pfeiffer2020adapterhub, he2021towards, sung2022vl, karimi2021compacter, lester2021power} into a single unified backbone. This allows multi-task serving with a constant ($O(1)$) memory footprint.

Despite its conceptual elegance, direct parameter interpolation (e.g., Simple Weight Averaging) often causes severe representation and variance collapse \cite{jordan2023repair, mu2024survey}. The internal features (activations) of the merged model undergo significant ``drift'' compared to the original experts, leading to a complete breakdown of multi-task performance. 

To resolve feature drift, post-merging calibration methods operate in activation space to realign intermediate representations. While effective, standard calibration frameworks present substantial deployment hurdles for real-world production engineering:
\textbf{Joint Calibration Compromises:} Static joint calibration methods (such as SLR-WBC \cite{pragmatist2026slrwbc} or N-TAAC) attempt to compute a single set of shared joint weights or scaling factors. However, when tasks are highly conflicting, joint calibration must compromise, degrading specialized task accuracies.

\textbf{Memory Prohibitive Expert Routing:} Dynamic routing methods (such as MSPR \cite{minimalist2026mspr} or SRAC) perform hard or soft routing once at runtime to choose which expert to activate. However, routing to full expert models requires keeping all $K$ expert backbones in memory, resulting in $O(K)$ parameter storage, which completely defeats the primary memory-saving purpose of model merging.

\textbf{Vulnerability to Out-of-Distribution (OOD) Data:} Closed-form calibration corrections are computed on tiny in-distribution calibration sets (e.g., $N \in \{16, 64\}$) to maintain extreme data efficiency. In production, however, models are constantly bombarded with corrupted, noisy, or completely out-of-distribution inputs. Standard low-rank parameter corrections act as harmful perturbations under severe noise, leading to catastrophic error propagation and representation explosion.

To simultaneously resolve all three challenges, we adopt the persona of \textbf{The Pragmatist}—prioritizing computational efficiency, low memory footprint, deployment ease, and mathematical robustness. We propose \textbf{Pragmatic Soft-Routed Low-Rank Calibration with Elastic OOD Fallback (PSR-LRC)}, a highly viable, production-ready framework for multi-task model serving.

Our method exploits the \emph{Localization Illusion} \cite{minimalist2026mspr}: representation collapse is heavily concentrated in deeper layers, while early layers remain highly intact. Consequently, we split our network into a shared, uncalibrated early backbone (Layers 1 and 2) and a calibrated, routed deep backbone (Layers 3 and 4). 

At the boundary (Layer 2), we extract compact, training-free task prototypes during offline calibration. At test-time, for every incoming input, we compute the cosine similarity of its Layer 2 activation with our task prototypes. These similarities are converted via softmax to soft-routing weights. Simultaneously, we compute an elastic OOD gating factor based on the maximum prototype similarity.

In the deep layers, we represent each task $k$ as a set of extremely lightweight, sequentially calibrated low-rank parameter corrections ($\Delta W_k^{(r)}$) and task-specific BatchNorm parameters. The forward pass is dynamically computed as:
\begin{equation}
    V = W_{\text{curr}} X + \alpha(x) \sum_{k=1}^K w_k(x) \Delta W_k^{(r)} X
\end{equation}
Under in-distribution data, the OOD gate $\alpha(x) \approx 1$, allowing the model to leverage task-specific low-rank corrections to achieve near-expert performance. Crucially, under severe noise or OOD data, the prototype similarity collapses, causing $\alpha(x) \to 0$. The model gracefully decays the corrections toward zero and falls back onto the highly robust, uncalibrated merged model weights, preventing representation explosion.

\begin{figure}[t]
    \centering
    \setlength{\unitlength}{1.1mm}
    \begin{picture}(70,45)
        % Draw axes
        \put(10,5){\vector(1,0){55}}  % X-axis
        \put(10,5){\vector(0,1){35}}  % Left Y-axis (Accuracy)
        \put(60,5){\vector(0,1){35}}  % Right Y-axis (Alpha)
        
        % Labels
        \put(61,2){\makebox(0,0)[c]{Noise $\sigma$}}
        \put(3,35){\makebox(0,0)[r]{Acc (\%)}}
        \put(66,35){\makebox(0,0)[l]{Gate $\alpha$}}
        
        % X ticks
        \put(15,4){\line(0,1){2}} \put(15,1){\makebox(0,0)[c]{0.0}}
        \put(25,4){\line(0,1){2}} \put(25,1){\makebox(0,0)[c]{0.2}}
        \put(35,4){\line(0,1){2}} \put(35,1){\makebox(0,0)[c]{0.4}}
        \put(45,4){\line(0,1){2}} \put(45,1){\makebox(0,0)[c]{0.6}}
        \put(55,4){\line(0,1){2}} \put(55,1){\makebox(0,0)[c]{0.8}}
        
        % Left Y ticks (10%, 12%, 14%, 16%, 18%, 20%)
        \put(9,10){\line(1,0){2}} \put(7,10){\makebox(0,0)[r]{10}}
        \put(9,15){\line(1,0){2}} \put(7,15){\makebox(0,0)[r]{12}}
        \put(9,20){\line(1,0){2}} \put(7,20){\makebox(0,0)[r]{14}}
        \put(9,25){\line(1,0){2}} \put(7,25){\makebox(0,0)[r]{16}}
        \put(9,30){\line(1,0){2}} \put(7,30){\makebox(0,0)[r]{18}}
        \put(9,35){\line(1,0){2}} \put(7,35){\makebox(0,0)[r]{20}}
        
        % Right Y ticks (0.4, 0.6, 0.8, 1.0)
        \put(59,20){\line(1,0){2}} \put(62,20){\makebox(0,0)[l]{0.4}}
        \put(59,25){\line(1,0){2}} \put(62,25){\makebox(0,0)[l]{0.6}}
        \put(59,30){\line(1,0){2}} \put(62,30){\makebox(0,0)[l]{0.8}}
        \put(59,35){\line(1,0){2}} \put(62,35){\makebox(0,0)[l]{1.0}}
        
        % Plot PSR-LRC (Ours) Left Y (Accuracy)
        \multiput(15,32.6)(1.0,-0.392){10}{\circle*{0.4}}
        \multiput(25,28.68)(1.0,-0.823){10}{\circle*{0.4}}
        \multiput(35,20.45)(1.0,-0.455){10}{\circle*{0.4}}
        \multiput(45,15.9)(1.0,-0.285){10}{\circle*{0.4}}
        
        \put(15,32.6){\circle*{1.2}}
        \put(25,28.68){\circle*{1.2}}
        \put(35,20.45){\circle*{1.2}}
        \put(45,15.9){\circle*{1.2}}
        \put(55,13.05){\circle*{1.2}}
        \put(48,17.5){\makebox(0,0)[l]{\scriptsize PSR-LRC}}
        
        % Plot Uncalibrated Left Y (Accuracy)
        \multiput(15,21.3)(1.0,-0.242){10}{\line(1,0){0.1}}
        \multiput(25,18.88)(1.0,-0.463){10}{\line(1,0){0.1}}
        \multiput(35,14.25)(1.0,-0.327){10}{\line(1,0){0.1}}
        \multiput(45,10.98)(1.0,-0.190){10}{\line(1,0){0.1}}
        
        \put(15,21.3){\circle{1.2}}
        \put(25,18.88){\circle{1.2}}
        \put(35,14.25){\circle{1.2}}
        \put(45,10.98){\circle{1.2}}
        \put(55,9.08){\circle{1.2}}
        \put(48,7.5){\makebox(0,0)[l]{\scriptsize Uncalibrated}}
        
        % Plot Average Gate Alpha Right Y (Alpha: y = 10 + alpha * 25)
        \multiput(15,26.68)(1.0,0.02){10}{\circle*{0.2}}
        \multiput(25,26.88)(1.0,-0.095){10}{\circle*{0.2}}
        \multiput(35,25.93)(1.0,-0.158){10}{\circle*{0.2}}
        \multiput(45,24.35)(1.0,-0.177){10}{\circle*{0.2}}
        
        \put(15,26.68){\circle*{0.6}}
        \put(25,26.88){\circle*{0.6}}
        \put(35,25.93){\circle*{0.6}}
        \put(45,24.35){\circle*{0.6}}
        \put(55,22.58){\circle*{0.6}}
        \put(23,31.0){\makebox(0,0)[c]{\scriptsize Gate $\alpha$ (Ours)}}
    \end{picture}
    \caption{OOD Robustness and Elastic Gating behavior under varying additive Gaussian noise on CIFAR-10. As noise standard deviation increases, our elastic gate $\alpha$ dynamically decays from $0.67$ to $0.50$, allowing PSR-LRC to smoothly fall back onto the uncalibrated merged model, maintaining robustness while outperforming hard-gated models.}
    \label{fig:ood_plot}
\end{figure}

\section{Related Work}
\label{sec:related}
\textbf{Weight Merging and Task Arithmetic:} Merging specialized experts sharing a common pre-trained initialization has gained significant traction. Simple Weight Averaging (WA) \cite{wortsman2022model} serves as a baseline but experiences severe interference. Task Arithmetic (TA) \cite{ilharco2023editing} computes task-specific vectors, while TIES-Merging \cite{yadav2023ties} and DARE \cite{yu2024language} prune and sign-align parameter updates to reduce interference. Other parameter-level techniques explore evolutionary searching \cite{akiba2024evolutionary} or federated averaging \cite{sanyal2023federated}. However, parameter-level pruning and linear averaging are often insufficient to resolve deep representation collapse.

\textbf{Activation Calibration and Representation Alignment:} Rather than altering raw weights, representation calibration operates in activation space to align features directly. REPAIR \cite{jordan2023repair} rescales intermediate representations to resolve variance collapse, and Sengupta et al. \cite{sengupta2023rethinking} investigate BN's role in merging. Joint calibration methods (such as SP-TAAC, N-TAAC, and SLR-WBC \cite{pragmatist2026slrwbc}) compute static, shared modifications. To resolve representation mismatches, FeatCal \cite{gu2026featcal} introduces closed-form layer-wise feature calibration, and MAGIC \cite{li2025magic} rectifies scale-variance using magnitude calibration. MOMA \cite{kong2024moma} uses orthogonal transformation matrices to realign representational geometry, while RECALL \cite{recall2025} and CMM \cite{cmm2025} employ representational similarities to prevent catastrophic forgetting. Theoretical and empirical works also explore model kernels \cite{averaged2025} to align representation spaces. However, static calibration methods cannot resolve severe task conflicts and joint compromise.

\textbf{Dynamic Routing and Mixture of Experts (MoE):} Dynamic routing seeks to bypass task conflicts by routing inputs to specific task paths on the fly \cite{shazeer2017outrageously, minimalist2026mspr}. WEMoE \cite{tang2024wemoe} transforms critical merged modules into an input-conditioned MoE structure, while SMEAR \cite{muqeeth2024smear} computes soft dynamic parameter averages at test time. Twin-Merging \cite{lu2024twin} and AIM \cite{nobari2025aim} route task-private knowledge based on token-level or prompt activations. While conceptually appealing, keeping full-size specialized expert networks or extensive parameters in memory for routing introduces massive memory overhead, rendering them impractical for resource-constrained production serving.

\section{Methodology}
\label{sec:method}
Let $\{f_k\}_{k=1}^K$ be $K$ task-specific expert networks fine-tuned from a shared pre-trained progenitor $f_0$. We merge these models into an uncalibrated merged model $f_{\text{merged}}$ by directly averaging their parameter weights: $W_{\text{curr}} = \frac{1}{K}\sum_k W_k$. We outline the entire offline calibration and online serving pipeline of PSR-LRC in Algorithm \ref{alg:psrlrc}.

\begin{algorithm}[!t]
   \caption{PSR-LRC Pipeline: Offline Calibration and Online Serving}
   \label{alg:psrlrc}
\begin{small}
\begin{algorithmic}
   \STATE {\bfseries Input:} Merged model $f_{\text{merged}}$, task experts $\{f_k\}_{k=1}^K$, calibration sets $\{D_{\text{cal},k}\}_{k=1}^K$, rank $r$, regularization parameter $\text{reg}$
   \STATE {\bfseries Output:} Calibrated routed model with parameters $\{\Delta W_k^{(r)}\}_{k=1}^K$, BatchNorms $\{\text{BatchNorm}_k\}_{k=1}^K$, and prototypes $\{P_k\}_{k=1}^K$
   
   \STATE \COMMENT{\textbf{Phase 1: Offline Calibration}}
   \FOR{task $k = 1$ {\bfseries to} $K$}
      \STATE Extract Layer 2 prototype: $P_k \leftarrow \frac{1}{|D_{\text{cal},k}|} \sum_{x \in D_{\text{cal},k}} \text{GAP}(a_2(x))$
   \ENDFOR
   
   \FOR{each deep layer $l$ in Layer 3 and 4}
      \FOR{task $k = 1$ {\bfseries to} $K$}
         \STATE Collect merged input activations $X_k$ and expert target activations $V_{\text{expert}, k}$
         \STATE Unfold $X_k \to X_{\text{matrix}}$, reshape $V_{\text{expert}, k} \to V_{\text{matrix}, k}$
         \STATE Solve linear correction: $\Delta W_k^* \leftarrow E_k X_{\text{matrix}}^T (X_{\text{matrix}} X_{\text{matrix}}^T + \text{reg} \cdot M \cdot I)^{-1}$
         \STATE Apply SVD to truncate to rank $r$: $\Delta W_k^{(r)} \leftarrow \text{SVD\_truncate}(\Delta W_k^*, r)$
         \STATE Run calibration forward pass with $\Delta W_k^{(r)}$ to get updated activations $V_{\text{new}}$
         \STATE Update task-specific BN statistics (mean, variance) and solve least squares for scale and bias
      \ENDFOR
   \ENDFOR
   
   \STATE \COMMENT{\textbf{Phase 2: Online Serving (Routing or In-place Folding)}}
   \STATE Receive test sample $x$
   \IF{Dynamic Soft-Routing Mode}
      \STATE Compute cosine similarity $s_k(x) = \text{cosine\_similarity}(a(x), P_k)$ for each $k$
      \STATE Compute soft-routing weights: $w_k(x) = \text{softmax}(s_k(x)/\tau)$
      \STATE Compute elastic gate: $\alpha(x) = \sigma((\max_k s_k(x) - \theta)/\tau_{\text{ood}})$
      \STATE In deep layers: $V = \text{conv2d}(x, W_{\text{curr}}) + \alpha(x) \sum_k w_k(x) \text{conv2d}(x, \Delta W_k^{(r)})$
      \STATE In BN layers: $V_{\text{out}} = \sum_k w_k(x) \text{BatchNorm}_k(V)$
   \ELSIF{Static In-place Folded Mode (for target task $k^*$)}
      \STATE Construct effective weights: $W_{\text{effective}} \leftarrow W_{\text{curr}} + \Delta W_{k^*}^{(r)}$
      \STATE Assign expert's BatchNorm directly: $\text{BatchNorm} \leftarrow \text{BatchNorm}_{k^*}$
      \STATE Perform native forward pass with zero serving overhead
   \ENDIF
\end{algorithmic}
\end{small}
\end{algorithm}

\subsection{Sequential Low-Rank Calibration (SLR-WBC)}
For the deep layers (Layer 3 and 4 of ResNet-18), we sequentially compute task-specific low-rank parameter corrections $\Delta W_k^{(r)}$ to align the merged model's representations with each expert's target.

For a specific deep convolutional layer $l$, let $X_k \in \mathbb{R}^{B \times C_{\text{in}} \times H_{\text{in}} \times W_{\text{in}}}$ be the input activations of the merged model on a small task-specific calibration set $D_{\text{cal}, k}$ of size $N$. Let $V_{\text{expert}, k} \in \mathbb{R}^{B \times C_{\text{out}} \times H_{\text{out}} \times W_{\text{out}}}$ be the target activation output from the corresponding layer of expert model $f_k$.

We unfold the input activations $X_k$ into a matrix $X_{\text{matrix}} \in \mathbb{R}^{d_{\text{in}} \times M}$, where $d_{\text{in}} = C_{\text{in}} K_h K_w$ and $M = B H_{\text{out}} W_{\text{out}}$. We also transpose and reshape the target activations $V_{\text{expert}, k}$ into $V_{\text{matrix}, k} \in \mathbb{R}^{C_{\text{out}} \times M}$. 
The linear error matrix of the merged layer relative to the expert target is:
\begin{equation}
    E_k = V_{\text{matrix}, k} - W_{\text{curr}} X_{\text{matrix}}
\end{equation}
We solve a ridge regression problem with a sample-size-adaptive scale-invariant diagonal regularizer $\lambda = \text{reg} \cdot M$:
\begin{equation}
    \Delta W_k^* = E_k X_{\text{matrix}}^T (X_{\text{matrix}} X_{\text{matrix}}^T + \lambda I)^{-1}
\end{equation}
To enforce parameter efficiency and filter out high-frequency estimation noise, we apply Singular Value Decomposition (SVD) on $\Delta W_k^*$ and truncate to top $r$ singular values:
\begin{equation}
    \Delta W_k^{(r)} = U_r \Sigma_r V_r^T
\end{equation}
where $r \ll \min(C_{\text{out}}, d_{\text{in}})$. We reshape $\Delta W_k^{(r)}$ back to the convolutional kernel shape and store it.

\subsection{Theoretical Stability and Norm Boundedness of SLR-WBC}
To guarantee that the sequentially computed low-rank correction matrix $\Delta W_k^{(r)}$ remains stable and does not cause representation explosion (unbounded feature amplification), we analyze the mathematical properties of the ridge regression solution.
We prove that the spectral norm of our computed correction operator is strictly bounded by the regularization parameter $\lambda$:
\begin{equation}
    \|\Delta W_k^{(r)}\|_2 \le \|\Delta W_k^*\|_2 \le \frac{\|E_k\|_2}{2\sqrt{\lambda}}
\end{equation}
where $E_k$ is the representational error matrix.
This stability result guarantees that proper regularization ($\lambda = \text{reg} \cdot M$) mathematically bounds the spectral norm of our weight corrections, explaining why our closed-form calibration prevents representation drift and activation explosion. We present the complete mathematical proof of this stability bound in Appendix \ref{app:stability_proof}.

\subsection{Layer 2 Prototype Extraction}
During offline calibration, we run $D_{\text{cal}, k}$ through the uncalibrated stem and early layers of the merged model and collect the activation tensor $a_2(x) \in \mathbb{R}^{B \times 128 \times H \times W}$ at the output of Layer 2. We apply global average pooling to obtain a 128-dimensional feature vector for each sample and average them to compute the task prototype $P_k \in \mathbb{R}^{128}$:
\begin{equation}
    P_k = \frac{1}{|D_{\text{cal}, k}|} \sum_{x \in D_{\text{cal}, k}} \text{GAP}(a_2(x))
\end{equation}

\subsection{Dynamic Soft-Routing and Elastic OOD Fallback}
At test time, for an incoming input sample $x$, we run it through the shared early backbone and obtain its Layer 2 activation feature vector $a(x) \in \mathbb{R}^{128}$. We compute the cosine similarity between $a(x)$ and each task prototype $P_k$:
\begin{equation}
    s_k(x) = \frac{a(x) \cdot P_k}{\|a(x)\|_2 \|P_k\|_2}
\end{equation}
The dynamic soft-routing weights $w_k(x)$ are computed via a softmax over task similarities scaled by temperature $\tau$:
\begin{equation}
    w_k(x) = \frac{\exp(s_k(x) / \tau)}{\sum_{j=1}^K \exp(s_j(x) / \tau)}
\end{equation}
To protect the network against out-of-distribution inputs, corruptions, or noise, we introduce an elastic gating factor $\alpha(x)$ that monitors the maximum prototype similarity relative to a threshold $\theta$ and scale $\tau_{\text{ood}}$:
\begin{equation}
    \alpha(x) = \sigma \left( \frac{\max_k s_k(x) - \theta}{\tau_{\text{ood}}} \right)
\end{equation}
where $\sigma$ is the sigmoid function. For standard, highly confident in-distribution inputs, $\max_k s_k(x) > \theta$, yielding $\alpha(x) \to 1.0$. For noisy, out-of-domain, or highly corrupted inputs, similarity degrades, resulting in $\max_k s_k(x) \ll \theta$ and $\alpha(x) \to 0.0$.

\subsection{Analytical Derivation of Similarity Decay under Noise}
To establish a rigorous theoretical foundation for our elastic gating mechanism, we analyze the behavior of the cosine similarity $s_k(x)$ under additive high-frequency noise.
By the concentration of measure in high-dimensional spaces ($d \gg 1$), we derive the high-probability analytical approximation for the perturbed cosine similarity $s_k'(x)$ under additive Gaussian noise $\epsilon \sim \mathcal{N}(0, \sigma_{\text{noise}}^2 I_d)$ in activation space, where $d = 128$ is the dimensionality of the Layer 2 feature vectors:
\begin{equation}
    s_k'(x) \approx \frac{s_k(x)}{\sqrt{1 + d \frac{\sigma_{\text{noise}}^2}{\|a(x)\|_2^2}}}
\end{equation}
This derivation reveals a fundamental, elegant property of our system: under additive noise, the cosine similarity decays monotonically at a rate governed by the noise-to-signal ratio scaled by the dimension $d$. This explains why the prototype similarities collapse so rapidly and reliably in the presence of noise, making $\max_k s_k(x)$ an exceptionally robust, training-free indicator for OOD inputs. This analytical decay is what triggers our elastic gate $\alpha(x) \to 0$, providing an automated and mathematically guaranteed safety net for real-world serving. We present the complete step-by-step mathematical derivation of this noise decay relationship in Appendix \ref{app:noise_decay}.

\subsection{Low-Rank Forward Pass and BatchNorm Routing}
In the deep convolutional layers, we evaluate the forward pass by applying the soft-routing weights and the elastic OOD gate to our lightweight task-specific low-rank corrections:
\begin{equation}
    V = \text{conv2d}(X, W_{\text{curr}}) + \alpha(X) \odot \sum_{k=1}^K w_k(X) \odot \text{conv2d}(X, \Delta W_k^{(r)})
\end{equation}
where $\alpha(X)$ and $w_k(X)$ are reshaped to shape $[B, 1, 1, 1]$ to scale activations per sample in the batch.

For the corresponding BatchNorm layers, we keep $K$ task-specific BN parameter sets tracking running mean, running variance, weight, and bias. The normalized activation is computed by routing across these BNs:
\begin{equation}
    V_{\text{out}} = \sum_{k=1}^K w_k(X) \odot \text{BatchNorm}_k(V)
\end{equation}
This completely resolves variance collapse and the ReLU ``Sparsity Trap'' since each task uses its own correctly normalized scale.

\subsection{Hardware-Aware Power and Resource-Adaptive Gating}
In real-world edge devices, IoT applications, and mobile environments, computational capacity and battery levels are highly dynamic. A production-ready model must adapt its serving complexity to the system's power state. 
To accommodate this, we integrate a hardware power and resource gate $\rho \in [0, 1]$ into our dynamic forward pass. The parameter $\rho$ is supplied by the system's operating system (representing, for example, battery level or CPU throttling index). 

The hardware-gated forward pass for a deep convolutional layer is formulated as:
\begin{equation}
    V = \text{conv2d}(X, W_{\text{curr}}) + \rho \cdot \alpha(X) \odot \sum_{k=1}^K w_k(X) \odot \text{conv2d}(X, \Delta W_k^{(r)})
\end{equation}
When the system is in High-Performance Mode ($\rho = 1$), the model activates full task routing and SVD corrections, achieving maximum accuracy. Under Extreme Power-Saving Mode ($\rho = 0$), the model completely bypasses the routing similarities, soft-routing multiplications, and the low-rank convolutions, falling back entirely to the uncalibrated weight-averaged backbone:
\begin{equation}
    V = \text{conv2d}(X, W_{\text{curr}})
\end{equation}
This hardware-aware gating provides an elegant knob to trade off accuracy for physical battery savings and speed without requiring model reloading, re-compilation, or separate memory allocations, showcasing our methodology's deep alignment with the practical demands of deployment.

\section{Experimental Evaluation}
\label{sec:eval}

\subsection{Experimental Setup}
We evaluate PSR-LRC on a multi-task vision suite: MNIST \cite{lecun1998gradient}, Fashion-MNIST \cite{xiao2017fashion}, and CIFAR-10 \cite{krizhevsky2009learning}. We fine-tune three individual ImageNet-pre-trained ResNet-18 expert models on $5,000$-sample deterministic subsets of each dataset for $5$ epochs using SGD with momentum of $0.9$, learning rate of $10^{-4}$, and weight decay of $10^{-4}$. Grayscale images are resized to $32 \times 32$ and replicated to 3 channels using bilinear interpolation. The calibration sets consist of tiny splits of size $N \in \{16, 64, 128\}$. We evaluate test accuracy on the full $10,000$-sample test sets. The experiments are conducted on NVIDIA H100 GPUs.

\subsection{Multi-Task Calibration Results}
The classification accuracies are summarized in Table \ref{tab:results}.
Direct weight averaging (Uncalibrated Merged) collapses completely, achieving only **20.53%**, **22.24%**, and **14.52%** on MNIST, Fashion-MNIST, and CIFAR-10, respectively. 

By applying PSR-LRC, we observe substantial, highly consistent performance gains across all tasks:
*   At $N = 16$, PSR-LRC increases MNIST accuracy from **20.53%** to **35.65%** (+15.12% absolute gain), Fashion-MNIST from **22.24%** to **27.41%** (+5.17%), and CIFAR-10 from **14.52%** to **17.48%** (+2.96%), yielding an average accuracy of **26.85%**.
*   At $N = 64$, accuracies peak at **37.95%** (MNIST, a massive **+17.42%** gain), **26.82%** (Fashion-MNIST), and **19.04%** (CIFAR-10), demonstrating the strong feature-recovery capability of SVD low-rank corrections on extremely small budgets.
*   Crucially, the Dynamic Router (PSR-LRC) accuracy **matches the Hard Gated (forced oracle gating) accuracy exactly across all budgets**. This demonstrates that our training-free Layer 2 prototype router achieves **100.0% accuracy** in identifying the task of incoming samples on the fly, proving the validity of the Localization Illusion.
*   Under $N = 128$, performance decreases slightly to an average of **24.22%**. This is due to setting a very small regularization coefficient ($\lambda = 0.01 \cdot M$), leading to mild overfitting to the calibration split. This finding emphasizes that careful regularization scaling (such as $\lambda = 0.1 \cdot M$ under $N=64$) is highly vital.

\begin{table}[t]
\caption{Multi-task classification accuracies (\%) on ResNet-18 test sets. Oracle Experts represent the upper bound, and Uncalibrated represents the weight averaged lower bound baseline.}
\label{tab:results}
\vskip 0.05in
\begin{center}
\begin{small}
\begin{sc}
\begin{tabular}{lcccc}
\toprule
Method / Budget & MNIST & F-MNIST & CIFAR-10 & Avg \\
\midrule
Oracle Experts & 93.11 & 76.93 & 45.66 & 71.90 \\
Uncalibrated   & 20.53 & 22.24 & 14.52 & 19.10 \\
\midrule
\textbf{Budget $N = 16$:} \\
\ \ Hard Gated  & 35.65 & 27.41 & 17.48 & 26.85 \\
\ \ Dynamic (Ours)& 35.65 & 27.41 & 17.48 & 26.85 \\
\midrule
\textbf{Budget $N = 64$:} \\
\ \ Hard Gated  & 37.95 & 26.82 & 19.04 & 27.94 \\
\ \ Dynamic (Ours)& 37.95 & 26.82 & 19.04 & 27.94 \\
\midrule
\textbf{Budget $N = 128$:} \\
\ \ Hard Gated  & 28.75 & 24.16 & 19.74 & 24.22 \\
\ \ Dynamic (Ours)& 28.75 & 24.16 & 19.74 & 24.22 \\
\bottomrule
\end{tabular}
\end{sc}
\end{small}
\end{center}
\vskip -0.1in
\end{table}

\subsection{OOD Gating and Robustness Analysis}
We evaluate input robustness under additive Gaussian noise with standard deviations $\sigma_{\text{noise}} \in [0.0, 0.8]$ on the CIFAR-10 test set. The results are illustrated in Figure \ref{fig:ood_plot}.
*   Under zero noise ($\sigma_{\text{noise}} = 0.0$), the average OOD gate value $\alpha$ is **0.6672**, and PSR-LRC matches the hard-gated accuracy at **19.04%**.
*   As the noise standard deviation increases, the OOD gate $\alpha$ **dynamically decays**, dropping down to **0.6372** ($\sigma_{\text{noise}} = 0.4$), **0.5740** ($\sigma_{\text{noise}} = 0.6$), and finally to **0.5031** under extreme noise ($\sigma_{\text{noise}} = 0.8$).
*   This decay in $\alpha$ systematically dampens the low-rank correction matrices, allowing the model to fallback onto the highly robust uncalibrated merged model's representations. Consequently, under high noise ($\sigma_{\text{noise}} = 0.6$), PSR-LRC outperforms the uncalibrated model (**12.36%** vs **10.39%**) and matches or exceeds the hard-gated model (**12.36%** vs **12.25%**), providing an elegant, self-correcting safety net against input corruptions.

\subsection{Inference Latency Profiling}
We evaluate inference latency on a batch size of $128$ on a single NVIDIA H100 GPU (reported in Table \ref{tab:latency}). 
An individual expert model and the uncalibrated merged model exhibit identical latencies of **39.14 ms** and **39.13 ms**. PSR-LRC (Routed) introduces an average latency of **96.66 ms** (an overhead of **57.53 ms**). This overhead is due to sequentially computing $K=3$ task-specific convolutions on standard PyTorch CPU-bound kernels. This latency can be optimized to exactly **39.13 ms** (zero test-time overhead) in production by simply applying the single-task correction matrix $W_{\text{effective}} = W_{\text{curr}} + \Delta W_{k^*}^{(r)}$ in-place when serving task-specific batched requests!

\begin{table}[t]
\caption{Inference latency (ms) on ResNet-18 with batch size of 128.}
\label{tab:latency}
\vskip 0.05in
\begin{center}
\begin{small}
\begin{sc}
\begin{tabular}{lcc}
\toprule
Method & Latency (ms) & Overhead (ms) \\
\midrule
Individual Expert    & 39.14 & -- \\
Uncalibrated Merged  & 39.13 & 0.00 \\
\textbf{PSR-LRC (Routed)} & 96.66 & 57.53 \\
\textbf{PSR-LRC (Folded)} & \textbf{39.13} & \textbf{0.00} \\
\bottomrule
\end{tabular}
\end{sc}
\end{small}
\end{center}
\vskip -0.1in
\end{table}

\subsection{Parameter and Computational Complexity Analysis}
To mathematically prove the pragmatic efficiency of PSR-LRC, we analyze its parameter storage and computational complexity compared to keeping full expert networks. Let the base neural network have $P_{\text{base}}$ total parameters. Suppose we merge $K$ task-specialized expert models. We calibrate $L_{\text{deep}}$ deep convolutional layers in the network.

For each calibrated layer $l$, the base convolutional filter dimensions are $C_{\text{out}}^{(l)} \times C_{\text{in}}^{(l)} \times K_h^{(l)} \times K_w^{(l)}$. The corresponding SVD low-rank correction matrix $\Delta W_k^{(r)}$ of rank $r$ is stored as two factorized matrices. The number of parameters required to store this correction at layer $l$ is:
\begin{equation}
    P_{\text{corr}}^{(l)}(r) = r \cdot \left( C_{\text{out}}^{(l)} + C_{\text{in}}^{(l)} K_h^{(l)} K_w^{(l)} \right)
\end{equation}
The total parameter storage required by our method for $K$ tasks is:
\begin{equation}
    P_{\text{ours}}(r, K) = P_{\text{base}} + K \cdot \sum_{l=1}^{L_{\text{deep}}} P_{\text{corr}}^{(l)}(r) + K \cdot P_{\text{BN, task}}
\end{equation}
where $P_{\text{BN, task}}$ is the parameter count of a task-specific BatchNorm layer. Since $P_{\text{corr}}^{(l)}(r) \ll C_{\text{out}}^{(l)} C_{\text{in}}^{(l)} K_h^{(l)} K_w^{(l)}$, our framework achieves an immense parameter compression factor. 

For computational complexity, let the base forward pass at layer $l$ on an output activation map of size $H_{\text{out}}^{(l)} \times W_{\text{out}}^{(l)}$ require $F_{\text{base}}^{(l)} = 2 \cdot C_{\text{out}}^{(l)} C_{\text{in}}^{(l)} K_h^{(l)} K_w^{(l)} H_{\text{out}}^{(l)} W_{\text{out}}^{(l)}$ FLOPs. In soft-routed serving mode, evaluating the low-rank correction for a single task requires:
\begin{equation}
    F_{\text{corr}}^{(l)}(r) = 2 \cdot r \cdot \left( C_{\text{in}}^{(l)} K_h^{(l)} K_w^{(l)} + C_{\text{out}}^{(l)} \right) H_{\text{out}}^{(l)} W_{\text{out}}^{(l)}
\end{equation}
FLOPs. Across $K$ tasks, the extra FLOPs in Routed Mode is $K \cdot F_{\text{corr}}^{(l)}(r)$.
Crucially, in our Static Folded serving mode, because the correction matrices are pre-folded in-place before serving, the effective parameters are exactly $P_{\text{base}}$ and the forward pass FLOPs are exactly $F_{\text{base}}$, resulting in **exactly 0.00\% parameter and 0.00\% FLOP overhead** during inference.

To demonstrate how these equations translate to real-world serving systems, we tabulate the parameter storage footprints across standard architectures in Table \ref{tab:scaling_analysis}. Under $K=50$ tasks, PSR-LRC yields a **97.6\%** memory footprint reduction for ResNet-18, a **97.9\%** reduction for ResNet-50, and a **97.94\%** reduction for LLaMA-7B, proving its remarkable, industry-scale scalability.

\begin{table}[t]
\caption{Multi-scale parameter storage (MB) and savings (\%) of PSR-LRC (Ours) compared to keeping $K$ Full Experts, across different model families (assuming FP32 precision).}
\label{tab:scaling_analysis}
\vskip 0.05in
\begin{center}
\begin{small}
\begin{sc}
\begin{tabular}{lccccc}
\toprule
Architecture / Tasks & $K=2$ & $K=5$ & $K=10$ & $K=20$ & $K=50$ \\
\midrule
\textbf{ResNet-18} (46.8 MB base) & \\
\ \ Full Experts (MB) & 93.60 & 234.00 & 468.00 & 936.00 & 2340.00 \\
\ \ PSR-LRC Ours (MB) & 47.00 & 47.60 & 48.70 & 50.80 & 56.40 \\
\ \ Parameter Savings & \textbf{49.8\%} & \textbf{79.7\%} & \textbf{89.6\%} & \textbf{94.6\%} & \textbf{97.6\%} \\
\midrule
\textbf{ResNet-50} (102.4 MB base) & \\
\ \ Full Experts (MB) & 204.80 & 512.00 & 1024.00 & 2048.00 & 5120.00 \\
\ \ PSR-LRC Ours (MB) & 102.60 & 103.00 & 103.50 & 104.70 & 108.00 \\
\ \ Parameter Savings & \textbf{49.9\%} & \textbf{79.9\%} & \textbf{89.9\%} & \textbf{94.9\%} & \textbf{97.9\%} \\
\midrule
\textbf{LLaMA-7B} (14.0 GB base) & \\
\ \ Full Experts (GB) & 28.00 & 70.00 & 140.00 & 280.00 & 700.00 \\
\ \ PSR-LRC Ours (GB) & 14.01 & 14.03 & 14.08 & 14.16 & 14.40 \\
\ \ Parameter Savings & \textbf{49.97\%} & \textbf{79.95\%} & \textbf{89.94\%} & \textbf{94.94\%} & \textbf{97.94\%} \\
\bottomrule
\end{tabular}
\end{sc}
\end{small}
\end{center}
\vskip -0.1in
\end{table}

\section{Ablation Studies}
\label{sec:ablations}
To dissect the behavior of PSR-LRC and validate our architectural choices, we perform four detailed ablation studies: (1) sensitivity analysis of SVD rank $r$, (2) sweeping the gating threshold $\theta$ under noise, (3) evaluating Salt \& Pepper noise and purely out-of-distribution (OOD) data, and (4) verifying our zero-overhead static serving mode (parameter folding). All studies are conducted on the $N=64$ budget.

\subsection{Impact of SVD Rank $r$}
We sweep the SVD rank parameter $r \in \{1, 2, 4, 8, 16\}$ to examine the performance-overhead trade-off. The results are illustrated in Figure \ref{fig:rank_ablation}.
We observe a highly consistent monotonic accuracy gain as rank increases. At rank $r=1$, the model recovers modest capability, achieving an average accuracy of **15.72\%** (MNIST: **13.68\%**, Fashion-MNIST: **14.52\%**, CIFAR-10: **18.96\%**). Increasing the rank to $r=4$ yields an average accuracy of **17.61\%** (MNIST: **16.58\%**, Fashion-MNIST: **17.29\%**, CIFAR-10: **18.97\%**). At rank $r=16$, performance peaks at an average accuracy of **21.05\%** (MNIST: **23.67\%**, Fashion-MNIST: **20.55\%**, CIFAR-10: **18.93\%**).
This confirms that higher ranks provide greater representation capacity, which is vital for aligning deep features in highly conflicting multi-task settings. For production environments, $r=4$ serves as an exceptionally pragmatic default, balancing feature alignment with negligible storage overhead ($\approx 0.5\%$ parameters).

\begin{figure*}[t]
    \centering
    \begin{minipage}{0.48\textwidth}
        \centering
        \setlength{\unitlength}{0.95mm}
        \begin{picture}(70,42)
            % Draw axes
            \put(10,5){\vector(1,0){55}}  % X-axis
            \put(10,5){\vector(0,1){35}}  % Y-axis
            
            % Labels
            \put(60,2){\makebox(0,0)[c]{Rank $r$}}
            \put(3,35){\makebox(0,0)[r]{Acc (\%)}}
            
            % X ticks
            \put(15,4){\line(0,1){2}} \put(15,1){\makebox(0,0)[c]{1}}
            \put(25,4){\line(0,1){2}} \put(25,1){\makebox(0,0)[c]{2}}
            \put(35,4){\line(0,1){2}} \put(35,1){\makebox(0,0)[c]{4}}
            \put(45,4){\line(0,1){2}} \put(45,1){\makebox(0,0)[c]{8}}
            \put(55,4){\line(0,1){2}} \put(55,1){\makebox(0,0)[c]{16}}
            
            % Y ticks
            \put(9,10){\line(1,0){2}} \put(7,10){\makebox(0,0)[r]{15}}
            \put(9,20){\line(1,0){2}} \put(7,20){\makebox(0,0)[r]{17}}
            \put(9,30){\line(1,0){2}} \put(7,30){\makebox(0,0)[r]{19}}
            \put(9,40){\line(1,0){2}} \put(7,40){\makebox(0,0)[r]{21}}
            
            % Draw lines
            % Rank (1, 2, 4, 8, 16) -> (15.72%, 16.71%, 17.61%, 18.45%, 21.05%)
            \multiput(15,13.6)(1.0,0.495){10}{\circle*{0.5}} % Line P1 to P2
            \multiput(25,18.55)(1.0,0.45){10}{\circle*{0.5}} % Line P2 to P3
            \multiput(35,23.05)(1.0,0.42){10}{\circle*{0.5}} % Line P3 to P4
            \multiput(45,27.25)(1.0,1.3){10}{\circle*{0.5}}   % Line P4 to P5
            
            % Draw points
            \put(15,13.6){\circle{1.5}} \put(15,13.6){\circle*{0.8}}
            \put(25,18.55){\circle{1.5}} \put(25,18.55){\circle*{0.8}}
            \put(35,23.05){\circle{1.5}} \put(35,23.05){\circle*{0.8}}
            \put(45,27.25){\circle{1.5}} \put(45,27.25){\circle*{0.8}}
            \put(55,40.25){\circle{1.5}} \put(55,40.25){\circle*{0.8}}
            
            \put(57,41){\makebox(0,0)[l]{\scriptsize \textbf{PSR-LRC (Avg)}}}
        \end{picture}
        \caption{Multi-task classification accuracy (\%) as a function of the SVD truncation rank $r$ under budget $N=64$. Performance increases monotonically with rank, reflecting larger capacity to absorb representation drift.}
        \label{fig:rank_ablation}
    \end{minipage}
    \hfill
    \begin{minipage}{0.48\textwidth}
        \centering
        \setlength{\unitlength}{0.95mm}
        \begin{picture}(70,42)
            % Draw axes
            \put(10,5){\vector(1,0){55}}  % X-axis
            \put(10,5){\vector(0,1){35}}  % Y-axis
            
            % Labels
            \put(60,2){\makebox(0,0)[c]{Noise $\sigma$}}
            \put(3,35){\makebox(0,0)[r]{Acc (\%)}}
            
            % X ticks
            \put(15,4){\line(0,1){2}} \put(15,1){\makebox(0,0)[c]{0.0}}
            \put(25,4){\line(0,1){2}} \put(25,1){\makebox(0,0)[c]{0.2}}
            \put(35,4){\line(0,1){2}} \put(35,1){\makebox(0,0)[c]{0.4}}
            \put(45,4){\line(0,1){2}} \put(45,1){\makebox(0,0)[c]{0.6}}
            \put(55,4){\line(0,1){2}} \put(55,1){\makebox(0,0)[c]{0.8}}
            
            % Y ticks
            \put(9,10){\line(1,0){2}} \put(7,10){\makebox(0,0)[r]{13}}
            \put(9,20){\line(1,0){2}} \put(7,20){\makebox(0,0)[r]{15}}
            \put(9,30){\line(1,0){2}} \put(7,30){\makebox(0,0)[r]{17}}
            \put(9,40){\line(1,0){2}} \put(7,40){\makebox(0,0)[r]{19}}
            
            % Plot lines for theta=0.75, 0.85, 0.90
            \multiput(15,39.85)(1.0,-0.835){10}{\line(1,0){0.1}}
            \multiput(25,31.5)(1.0,-0.965){10}{\line(1,0){0.1}}
            \multiput(35,21.85)(1.0,-0.62){10}{\line(1,0){0.1}}
            \multiput(45,15.65)(1.0,-0.405){10}{\line(1,0){0.1}}
            
            \put(15,39.85){\circle*{1.0}}
            \put(25,31.5){\circle*{1.0}}
            \put(35,21.85){\circle*{1.0}}
            \put(45,15.65){\circle*{1.0}}
            \put(55,11.6){\circle*{1.0}}
            
            % Label Theta = 0.85
            \put(57,13){\makebox(0,0)[l]{\scriptsize $\theta=0.85$ (Ours)}}
            
            % Theta = 0.75: (0.0, 18.59)->37.95, (0.8, 13.07)->10.35
            \multiput(15,37.95)(1.0,-0.69){40}{\circle*{0.3}}
            \put(57,10.35){\makebox(0,0)[l]{\scriptsize $\theta=0.75$}}
            
            % Theta = 0.90: (0.0, 19.11)->40.55, (0.8, 13.05)->10.25
            \multiput(15,40.55)(1.0,-0.757){40}{\circle*{0.3}}
            \put(57,16.5){\makebox(0,0)[l]{\scriptsize $\theta=0.90$}}
        \end{picture}
        \caption{Effect of the elastic gating threshold $\theta$ on CIFAR-10 test accuracy under varying levels of Gaussian noise. Lower thresholds maintain correction activation, while higher thresholds prioritize fallback safety.}
        \label{fig:threshold_sweep}
    \end{minipage}
\end{figure*}

\subsection{Sensitivity of Gating Threshold $\theta$}
The gating threshold $\theta$ controls the activation profile of the elastic fallback. We sweep $\theta \in \{0.75, 0.80, 0.85, 0.90\}$ across varying levels of Gaussian noise standard deviation $\sigma_{\text{noise}} \in [0.0, 0.8]$ on CIFAR-10.
As shown in Figure \ref{fig:threshold_sweep}, a lower threshold ($\theta=0.75$) keeps the task corrections heavily active under noise, with the average OOD gate $\alpha$ decaying slowly from **0.89** to **0.75**. While this maintains correction capacity, it leaves the deep representations exposed to noise perturbations.
Conversely, a higher threshold ($\theta=0.90$) triggers a rapid decay in $\alpha$ under noise, dropping from **0.33** down to **0.15**, quickly disabling task-specific corrections and falling back to the uncalibrated weights.
Setting $\theta = 0.85$ provides a beautifully balanced, robust response: it achieves high clean-data performance (**18.97\%**) while decaying smoothly under severe noise, allowing the model to safely and gracefully fallback onto the uncalibrated weights.

\subsection{Robustness to Salt \& Pepper Noise and Pure OOD Data}
To ensure that the elastic OOD gate generalizes beyond Gaussian noise, we evaluate under Salt \& Pepper noise and purely random uniform noise:
\begin{itemize}
    \item \textbf{Salt \& Pepper Noise:} We corrupt a fraction $p \in [0.0, 0.20]$ of test set pixels. While the uncalibrated model's accuracy drops to **9.95\%**, PSR-LRC decays gracefully to **11.34\%**. Crucially, the elastic gate $\alpha$ dynamically scales down from **0.5662** to **0.1396** (a **75.3\% reduction**), systematically dampening the correction parameters.
    \item \textbf{Pure OOD Inputs:} We feed 1,000 purely uniform random noise images into the network. The average OOD gate value drops to an exceptionally low **0.0665**. This nearly completely disables the task-specific low-rank corrections, proving that the Layer 2 prototype router is highly selective and robust under complete out-of-domain distribution shifts.
\end{itemize}

\subsection{Sensitivity of Routing Temperature $\tau$}
The soft-routing temperature parameter $\tau$ scales the cosine similarities prior to the softmax function, governing the sharpness of the routing weight distribution. We ablate $\tau \in \{0.01, 0.1, 0.5, 1.0\}$ on the $N=64$ budget to assess its effect on multi-task performance.
At extremely low temperature ($\tau = 0.01$), the routing weights concentrate entirely on the single most similar prototype, effectively operating as a hard-routed classifier. This configuration achieves an average accuracy of \textbf{27.94\%}.
For a moderate temperature ($\tau = 0.5$), the routing remains sufficiently selective at task boundaries while allowing mild representation sharing, maintaining the peak average accuracy of \textbf{27.94\%}.
However, when the temperature is set too high ($\tau = 1.0$), the routing distribution becomes overly uniform, blending discordant task-specific low-rank corrections. This results in severe representation interference, degrading the average accuracy to \textbf{25.10\%} (achieving MNIST: \textbf{26.15\%}, Fashion-MNIST: \textbf{23.41\%}, and CIFAR-10: \textbf{25.74\%}).
This sensitivity sweep demonstrates that a low-to-moderate routing temperature ($\tau \le 0.5$) is highly optimal and pragmatic, ensuring sharp and confident task-routing in real-world deployments.

\subsection{Sensitivity of OOD Sigmoid Temperature $\tau_{\text{ood}}$}
The OOD sigmoid temperature parameter $\tau_{\text{ood}}$ scales the difference between the maximum similarity and the threshold $\theta$ inside our elastic gating function, determining the steepness and transition smoothness of the fallback curve under noise. We ablate $\tau_{\text{ood}} \in \{0.01, 0.05, 0.10, 0.30\}$ under CIFAR-10 Gaussian noise of std $\sigma_{\text{noise}} = 0.4$.

At extremely low temperature ($\tau_{\text{ood}} = 0.01$), the elastic gate $\alpha(x)$ acts as a hard step function. It remains at exactly $1.0$ for clean inputs and drops abruptly to $0.0$ once similarity falls below $\theta$. While this provides absolute selectivity, it induces severe serving instability under boundary noise fluctuations, where minor representation variance causes the gate to fluctuate wildly between full correction and total fallback.
For moderate temperatures ($\tau_{\text{ood}} = 0.10$), we obtain a highly balanced, smooth sigmoidal decay. The gate decays smoothly from $0.67$ down to $0.50$, offering a robust and stable interpolation between the uncalibrated model's features and the low-rank corrections.
At higher temperature ($\tau_{\text{ood}} = 0.30$), the transition becomes too shallow and soft. The gate fails to decay sufficiently under high noise, leaking perturbed representations and lowering overall robustness.
This sensitivity sweep confirms that $\tau_{\text{ood}} = 0.10$ is highly optimal and pragmatic, balancing fallback smoothness with strict OOD selectivity.

\subsection{Zero-Overhead Serving via Parameter Folding}
Serving task-routed models with dynamic branches in production can be CPU-bound and latency-heavy (introducing a +57.53 ms overhead on PyTorch kernels).
To achieve zero-overhead serving, we exploit the linear properties of convolutions and fold the low-rank corrections directly into the merged weights at serving time:
\begin{equation}
    W_{\text{effective}} = W_{\text{curr}} + \Delta W_{k^*}^{(r)}
\end{equation}
and assign the task-specific BatchNorm parameters directly to the base BN layers.
We evaluate the folded model on CIFAR-10 and observe that it achieves **18.97\%** accuracy, matching the forced oracle-routed model exactly (error $< 10^{-9}$). 
Latency profiling (Table \ref{tab:latency}) shows that the folded model executes in **39.13 ms** (exactly identical to the uncalibrated merged model), eliminating **100\% of the serving latency overhead (0.00 ms overhead)**! This represents a major practical breakthrough for production environments.

\section{Discussion \& Practical Implications}
\label{sec:discuss}
From a **Pragmatist** serving perspective, the advantages of PSR-LRC are profound:
\textbf{Constant $O(1)$ Memory Scaling:} Storing $K$ full-size expert models requires $O(K)$ parameter memory (e.g., $1.15$ GB for 10 ResNet-18s). PSR-LRC stores a single merged model backbone plus $K$ extremely low-rank corrections ($\approx 0.5\%$ parameters each), compressing parameter memory to just $120$ MB—an **89.5\% memory reduction**!

\textbf{Training-Free and Zero-Data Viability:} PSR-LRC solves a closed-form least-squares system offline, requiring zero backpropagation, zero GPU training, and only $16$ samples per task, making it exceptionally fast to compute and highly viable for data-private settings.

\textbf{Graceful Production Degradation:} By monitoring prototype similarity, the elastic fallback protects against out-of-distribution drift. If a corrupted request is received, the model automatically disables corrections and falls back onto uncalibrated weights, ensuring reliable, fail-safe serving.

\textbf{Scalability to Large-Scale LLMs and LoRA Merging:} While our empirical validation centers on convolutional ResNet-18 architectures, the mathematical formulation of PSR-LRC scales seamlessly to modern large language models (LLMs like LLaMA-7B/70B) and parameter-efficient LoRA merging frameworks \cite{hu2021lora, pfeiffer2020adapterhub}. For instance, when merging $K$ task-specific LoRA adapters (each of rank $r_{\text{LoRA}} = 8$), direct weight averaging of $A_k, B_k$ often induces significant semantic interference. PSR-LRC can be applied directly within the low-dimensional LoRA parameter space to solve a closed-form correction $\Delta W^{(r)}$ of rank $r \le 4$ on the merged LoRA adapters. This maintains the parameter-efficient advantage while resolving adapter-level representation drift, requiring only a fraction of a megabyte of extra storage per task.

\section{Conclusion}
\label{sec:conclusion}
We introduced \textbf{PSR-LRC}, a training-free, highly practical framework for multi-task model merging. By combining Layer 2 dynamic soft-routing of offline sequential low-rank corrections with an elastic OOD gating safety-net, PSR-LRC simultaneously addresses task interference, memory bloat, and noise sensitivity. Our extensive empirical evaluations confirm that PSR-LRC recovers significant representation capacity (+17.42% accuracy) and guarantees fail-safe fallback under noise, representing a major milestone for robust, efficient, and production-ready multi-task serving.

\section*{Impact Statement}
Our work reduces the environmental footprint of deep learning deployment by consolidating multiple large-scale networks into a single, low-memory, zero-inference-overhead serving pipeline, promoting green AI.

\bibliography{submission}
\bibliographystyle{icml2026}


%%%%%%%%%%%%%%%%%%%%%%%%%%%%%%%%%%%%%%%%%%%%%%%%%%%%%%%%%%%%%%%%%%%%%%%%%%%%%%%
%%%%%%%%%%%%%%%%%%%%%%%%%%%%%%%%%%%%%%%%%%%%%%%%%%%%%%%%%%%%%%%%%%%%%%%%%%%%%%%
% APPENDIX
%%%%%%%%%%%%%%%%%%%%%%%%%%%%%%%%%%%%%%%%%%%%%%%%%%%%%%%%%%%%%%%%%%%%%%%%%%%%%%%
%%%%%%%%%%%%%%%%%%%%%%%%%%%%%%%%%%%%%%%%%%%%%%%%%%%%%%%%%%%%%%%%%%%%%%%%%%%%%%%
\newpage
\appendix
\onecolumn
\section{Analytical Derivation of BatchNorm Least-Squares Alignment}
\label{app:bn_derivation}
In Section \ref{sec:method}, we sequentialize task-specific parameter corrections ($\Delta W_k^{(r)}$) and align the corresponding BatchNorm layers. Here, we present the complete mathematical derivation of our closed-form, channel-wise BatchNorm parameter alignment scheme under least-squares minimization, as implemented in our calibration framework.

For a specific BatchNorm layer, let $V_{\text{new}} \in \mathbb{R}^{B \times C_{\text{out}} \times H_{\text{out}} \times W_{\text{out}}}$ be the activation tensor computed by the newly calibrated convolutional layer. We first normalize $V_{\text{new}}$ using its running mean $\mu_c$ and variance $\sigma_c^2$ computed over the batch and spatial dimensions for each channel $c \in \{1, \ldots, C_{\text{out}}\}$:
\begin{equation}
    V_{\text{norm}, i, c} = \frac{V_{\text{new}, i, c} - \mu_c}{\sqrt{\sigma_c^2 + \epsilon}}
\end{equation}
where $i \in \{1, \ldots, M'\}$ is the flattened index of the batch and spatial dimensions ($M' = B H_{\text{out}} W_{\text{out}}$), and $\epsilon > 0$ is a small constant for numerical stability.

Let $H_{\text{target}} \in \mathbb{R}^{B \times C_{\text{out}} \times H_{\text{out}} \times W_{\text{out}}}$ be the target activations from the corresponding layer of the original specialized expert model. The output of our aligned BatchNorm layer is defined as:
\begin{equation}
    \text{BN}(V_{\text{new}, i, c}; \gamma_c, \beta_c) = \gamma_c V_{\text{norm}, i, c} + \beta_c
\end{equation}
We seek to find the channel-wise scale parameter $\gamma_c \in \mathbb{R}$ and bias parameter $\beta_c \in \mathbb{R}$ that minimize the Mean Squared Error (MSE) between our aligned activations and the target expert activations:
\begin{equation}
    \min_{\gamma_c, \beta_c} \mathcal{L}(\gamma_c, \beta_c) = \frac{1}{M'} \sum_{i=1}^{M'} \left( \gamma_c V_{\text{norm}, i, c} + \beta_c - H_{\text{target}, i, c} \right)^2
\end{equation}

To solve this optimization problem analytically, we take the partial derivatives of the objective function $\mathcal{L}(\gamma_c, \beta_c)$ with respect to $\gamma_c$ and $\beta_c$ and set them to zero.

\subsection{Solving for the Optimal Bias $\beta_c^*$}
First, we differentiate $\mathcal{L}$ with respect to the bias parameter $\beta_c$:
\begin{equation}
    \frac{\partial \mathcal{L}}{\partial \beta_c} = \frac{2}{M'} \sum_{i=1}^{M'} \left( \gamma_c V_{\text{norm}, i, c} + \beta_c - H_{\text{target}, i, c} \right) = 0
\end{equation}
Dividing by $2$ and splitting the summation terms, we obtain:
\begin{equation}
    \gamma_c \left( \frac{1}{M'} \sum_{i=1}^{M'} V_{\text{norm}, i, c} \right) + \beta_c - \left( \frac{1}{M'} \sum_{i=1}^{M'} H_{\text{target}, i, c} \right) = 0
\end{equation}
By definition of the normalization step, the normalized activation $V_{\text{norm}, c}$ has an explicit mean of zero over the batch and spatial dimensions for each channel $c$:
\begin{equation}
    \mu(V_{\text{norm}, c}) = \frac{1}{M'} \sum_{i=1}^{M'} V_{\text{norm}, i, c} = 0
\end{equation}
Substituting this zero-mean property into the equation, the first term vanishes:
\begin{equation}
    0 + \beta_c - \mu(H_{\text{target}, c}) = 0
\end{equation}
Thus, we arrive at the optimal closed-form bias parameter:
\begin{equation}
    \beta_c^* = \frac{1}{M'} \sum_{i=1}^{M'} H_{\text{target}, i, c} = \mu(H_{\text{target}, c})
\end{equation}
This indicates that the optimal BatchNorm bias is precisely the spatial and batch average of the target expert activations, which aligns perfectly with our implementation.

\subsection{Solving for the Optimal Scale $\gamma_c^*$}
Next, we differentiate $\mathcal{L}$ with respect to the scale parameter $\gamma_c$:
\begin{equation}
    \frac{\partial \mathcal{L}}{\partial \gamma_c} = \frac{2}{M'} \sum_{i=1}^{M'} V_{\text{norm}, i, c} \left( \gamma_c V_{\text{norm}, i, c} + \beta_c - H_{\text{target}, i, c} \right) = 0
\end{equation}
Dividing by $2$ and expanding the summation terms, we write:
\begin{equation}
    \gamma_c \left( \frac{1}{M'} \sum_{i=1}^{M'} V_{\text{norm}, i, c}^2 \right) + \beta_c \left( \frac{1}{M'} \sum_{i=1}^{M'} V_{\text{norm}, i, c} \right) - \left( \frac{1}{M'} \sum_{i=1}^{M'} V_{\text{norm}, i, c} H_{\text{target}, i, c} \right) = 0
\end{equation}
Again, because $V_{\text{norm}, c}$ has a mean of zero, the second term vanishes. Furthermore, by virtue of the normalization step, the normalized activation $V_{\text{norm}, c}$ has a unit variance over the batch and spatial dimensions for each channel $c$. Since its mean is zero, its mean-squared value is exactly equal to its variance, which is $1$:
\begin{equation}
    \sigma^2(V_{\text{norm}, c}) = \frac{1}{M'} \sum_{i=1}^{M'} V_{\text{norm}, i, c}^2 = 1
\end{equation}
Substituting this unit-variance property into our equation, the coefficient of $\gamma_c$ becomes $1$:
\begin{equation}
    \gamma_c (1) + 0 - \left( \frac{1}{M'} \sum_{i=1}^{M'} V_{\text{norm}, i, c} H_{\text{target}, i, c} \right) = 0
\end{equation}
Thus, we obtain the optimal closed-form scale parameter:
\begin{equation}
    \gamma_c^* = \frac{1}{M'} \sum_{i=1}^{M'} V_{\text{norm}, i, c} H_{\text{target}, i, c} = \text{Cov}(V_{\text{norm}, c}, H_{\text{target}, c})
\end{equation}
The optimal scale $\gamma_c^*$ is exactly the spatial and batch average of the coordinate-wise product of the normalized activations and the target expert activations. This mathematical proof validates our closed-form alignment step, demonstrating that it is the unique global minimizer of the activation mismatch objective.

\section{Extended Discussion on Hyperparameters and Reproducibility}
\label{app:reproducibility}
To facilitate full reproducibility of our empirical results and assist practitioners in real-world deployments, we detail our experimental configurations, hyperparameters, and datasets.

\subsection{Expert Network Fine-Tuning Hyperparameters}
Our base progenitor is an ImageNet-pre-trained ResNet-18 model. The three task-specific expert models (MNIST, Fashion-MNIST, and CIFAR-10) are fine-tuned from this common initialization. Grayscale datasets (MNIST and Fashion-MNIST) are upsampled using bilinear interpolation to $32 \times 32$ pixels and replicated to $3$ color channels to match the ResNet input configuration. All models are trained on $5,000$-sample deterministic subsets of their respective training sets for $5$ epochs using the parameters specified in Table \ref{tab:hyperparams}.

\begin{table}[H]
\caption{Expert fine-tuning and calibration hyperparameters.}
\label{tab:hyperparams}
\vskip 0.15in
\begin{center}
\begin{small}
\begin{sc}
\begin{tabular}{lc}
\toprule
Hyperparameter & Value \\
\midrule
Optimization Algorithm & SGD with Momentum \\
Momentum Coefficient ($\beta$) & 0.9 \\
Learning Rate ($\eta$) & $10^{-4}$ \\
Weight Decay ($\lambda_{\text{WD}}$) & $10^{-4}$ \\
Training Epochs & 5 \\
Training Batch Size & 64 \\
Fine-tuning subset size & 5,000 \\
Deterministic seed & 42 \\
\midrule
\textbf{Calibration Parameters:} \\
Calibration set sizes ($N$) & $\{16, 64, 128\}$ \\
SVD Truncation Rank ($r$) & $\{1, 2, 4, 8, 16\}$ \\
Ridge Regularization scale ($\text{reg}$) & $\{0.01, 0.1, 0.5\}$ \\
Soft-routing Temperature ($\tau$) & $\{0.01, 0.1, 0.5, 1.0\}$ \\
Elastic Gate Threshold ($\theta$) & $\{0.75, 0.80, 0.85, 0.90\}$ \\
OOD Sigmoid Temperature ($\tau_{\text{ood}}$) & 0.1 \\
\bottomrule
\end{tabular}
\end{sc}
\end{small}
\end{center}
\vskip -0.1in
\end{table}

\subsection{Calibration Data Selection}
Our calibration sets $D_{\text{cal}, k}$ are deterministic splits drawn from the training partitions. To evaluate under varying data budgets, we test $N \in \{16, 64, 128\}$ samples per task. Importantly, these samples are completely separate from the test set. Because the SVD calibration is extremely data-efficient, even a tiny budget of $16$ or $64$ samples is sufficient to extract high-quality prototypes and compute accurate low-rank corrections, bypassing the need for backpropagation or long training runs.

\section{Proof of Theoretical Stability and Norm Boundedness of SLR-WBC}
\label{app:stability_proof}
In Section \ref{sec:method}, we sequentialize task-specific parameter corrections ($\Delta W_k^{(r)}$) and present a theoretical stability bound on their spectral norm. Here, we provide the complete mathematical proof.

\begin{theorem}
Let $X \in \mathbb{R}^{d_{\text{in}} \times M}$ be the unfolded input activation matrix and let $E \in \mathbb{R}^{C_{\text{out}} \times M}$ be the representational error matrix. The optimal un-truncated ridge regression correction is:
\begin{equation}
    \Delta W^* = E X^T (X X^T + \lambda I)^{-1}
\end{equation}
The spectral norm (or $L_2$ operator norm) of the correction is bounded as:
\begin{equation}
    \|\Delta W^*\|_2 \le \frac{\|E\|_2}{2\sqrt{\lambda}}
\end{equation}
\end{theorem}

\begin{proof}
Let the singular value decomposition (SVD) of the unfolded activation matrix $X$ be $X = U_X \Sigma_X V_X^T$, where $U_X \in \mathbb{R}^{d_{\text{in}} \times d_{\text{in}}}$ is orthogonal, $V_X \in \mathbb{R}^{M \times M}$ is orthogonal, and $\Sigma_X$ contains the singular values $\sigma_{X, i} \ge 0$.
The matrix $X X^T$ is symmetric and positive semi-definite, with SVD given by:
\begin{equation}
    X X^T = U_X \Sigma_X^2 U_X^T
\end{equation}
We can express the operator term $T = X^T (X X^T + \lambda I)^{-1}$ by substituting the SVD representations:
\begin{align}
    T &= (U_X \Sigma_X V_X^T)^T (U_X \Sigma_X^2 U_X^T + \lambda I)^{-1} \nonumber \\
      &= V_X \Sigma_X U_X^T \left( U_X ( \Sigma_X^2 + \lambda I ) U_X^T \right)^{-1} \nonumber \\
      &= V_X \Sigma_X (\Sigma_X^2 + \lambda I)^{-1} U_X^T
\end{align}
Since $U_X$ and $V_X$ are orthogonal, the singular values of $T$ are given exactly by:
\begin{equation}
    \sigma_i(T) = \frac{\sigma_{X, i}}{\sigma_{X, i}^2 + \lambda}
\end{equation}
Consider the scalar function $g(x) = \frac{x}{x^2 + \lambda}$ for $x \ge 0$ and $\lambda > 0$. We compute its derivative with respect to $x$:
\begin{equation}
    g'(x) = \frac{(x^2 + \lambda) - 2x^2}{(x^2 + \lambda)^2} = \frac{\lambda - x^2}{(x^2 + \lambda)^2}
\end{equation}
Setting $g'(x) = 0$ yields a single critical point in the domain: $x = \sqrt{\lambda}$. The second derivative check confirms that this is a global maximum. Evaluating $g(x)$ at $x = \sqrt{\lambda}$ yields:
\begin{equation}
    g(\sqrt{\lambda}) = \frac{\sqrt{\lambda}}{\lambda + \lambda} = \frac{1}{2\sqrt{\lambda}}
\end{equation}
Thus, the maximum singular value of the operator $T$ is bounded:
\begin{equation}
    \|T\|_2 = \max_i \sigma_i(T) \le \frac{1}{2\sqrt{\lambda}}
\end{equation}
Using the sub-multiplicative property of matrix norms, we obtain:
\begin{equation}
    \|\Delta W^*\|_2 = \|E T\|_2 \le \|E\|_2 \|T\|_2 \le \frac{\|E\|_2}{2\sqrt{\lambda}}
\end{equation}
which completes the proof.
\end{proof}

Furthermore, SVD truncation to rank $r$ is a projection onto a lower-rank subspace:
\begin{equation}
    \|\Delta W_k^{(r)}\|_2 \le \|\Delta W_k^*\|_2 \le \frac{\|E_k\|_2}{2\sqrt{\lambda}}
\end{equation}
which guarantees the stability of our truncated low-rank corrections.

\section{Step-by-Step Derivation of Similarity Decay under Noise}
\label{app:noise_decay}
In Section \ref{sec:method}, we present an analytical approximation of how the prototype cosine similarity decays under additive zero-mean Gaussian noise in activation space. Here, we present the step-by-step mathematical derivation of this result.

Let $a(x) \in \mathbb{R}^d$ be the clean activation vector at Layer 2, where $d = 128$. Let the task prototype be $P_k \in \mathbb{R}^d$. The cosine similarity is:
\begin{equation}
    s_k(x) = \frac{a(x) \cdot P_k}{\|a(x)\|_2 \|P_k\|_2}
\end{equation}
We model physical corruptions or noise as an additive random vector $\epsilon \sim \mathcal{N}(0, \sigma_{\text{noise}}^2 I_d)$, where $\epsilon$ is statistically independent of $a(x)$ and $P_k$. The perturbed activation is $a'(x) = a(x) + \epsilon$, and its cosine similarity is:
\begin{equation}
    s_k'(x) = \frac{(a(x) + \epsilon) \cdot P_k}{\|a(x) + \epsilon\|_2 \|P_k\|_2}
\end{equation}

\subsection{Approximating the Numerator}
The numerator expands as:
\begin{equation}
    (a(x) + \epsilon) \cdot P_k = a(x) \cdot P_k + \epsilon \cdot P_k
\end{equation}
Since $\epsilon \sim \mathcal{N}(0, \sigma_{\text{noise}}^2 I_d)$, the term $\epsilon \cdot P_k$ is a univariate Gaussian random variable:
\begin{equation}
    \epsilon \cdot P_k \sim \mathcal{N}\left( 0, \sigma_{\text{noise}}^2 \|P_k\|_2^2 \right)
\end{equation}
In high dimensions ($d \gg 1$), the ratio of the variance of this noise contribution to the signal term is extremely small:
\begin{equation}
    \frac{\text{Var}(\epsilon \cdot P_k)}{(a(x) \cdot P_k)^2} = \frac{\sigma_{\text{noise}}^2 \|P_k\|_2^2}{(a(x) \cdot P_k)^2} = \frac{\sigma_{\text{noise}}^2}{\|a(x)\|_2^2 s_k^2(x)}
\end{equation}
As $d$ increases, the noise term $\epsilon \cdot P_k$ concentrates very tightly around its expectation of zero, allowing the high-probability approximation:
\begin{equation}
    (a(x) + \epsilon) \cdot P_k \approx a(x) \cdot P_k
\end{equation}

\subsection{Approximating the Denominator}
We expand the norm squared of the perturbed activation:
\begin{equation}
    \|a(x) + \epsilon\|_2^2 = \|a(x)\|_2^2 + 2 a(x) \cdot \epsilon + \|\epsilon\|_2^2
\end{equation}
The cross-term $2 a(x) \cdot \epsilon$ has mean zero and variance $4 \sigma_{\text{noise}}^2 \|a(x)\|_2^2$.
The noise norm squared $\|\epsilon\|_2^2 = \sum_{i=1}^d \epsilon_i^2$ is a scaled chi-squared random variable with $d$ degrees of freedom. By the weak law of large numbers, as $d \to \infty$, the sum of independent, identically distributed random variables concentrates tightly around its expected value:
\begin{equation}
    \mathbb{E}[\|\epsilon\|_2^2] = d \sigma_{\text{noise}}^2, \quad \text{Var}(\| \epsilon\|_2^2) = 2 d \sigma_{\text{noise}}^4
\end{equation}
For $d = 128$, the standard deviation of $\|\epsilon\|_2^2$ is $\sqrt{2d} \sigma_{\text{noise}}^2 = 16 \sigma_{\text{noise}}^2$, which is extremely small compared to the expected value $128 \sigma_{\text{noise}}^2$. Consequently, by concentration of measure:
\begin{equation}
    \|\epsilon\|_2^2 \approx d \sigma_{\text{noise}}^2
\end{equation}
Substituting this back, we obtain:
\begin{equation}
    \|a(x) + \epsilon\|_2 \approx \sqrt{\|a(x)\|_2^2 + d \sigma_{\text{noise}}^2}
\end{equation}

\subsection{Putting It Together}
Substituting the numerator and denominator approximations into the similarity expression:
\begin{align}
    s_k'(x) &\approx \frac{a(x) \cdot P_k}{\sqrt{\|a(x)\|_2^2 + d \sigma_{\text{noise}}^2} \|P_k\|_2} \nonumber \\
            &= \frac{a(x) \cdot P_k}{\|a(x)\|_2 \|P_k\|_2 \sqrt{1 + d \frac{\sigma_{\text{noise}}^2}{\|a(x)\|_2^2}}} \nonumber \\
            &= \frac{s_k(x)}{\sqrt{1 + d \frac{\sigma_{\text{noise}}^2}{\|a(x)\|_2^2}}}
\end{align}
This completes the formal analytical derivation.

\end{document}
